\section{The hypothesis}

\begin{hypothesis}\label{hyp:ca}
For $\alpha_*\in(0,1]$ and a function $\Theta(n,M,l)$, $\mathrm H(\alpha_*,\Theta)$ is the
following assertion. For every finite field $\E\supseteq\F$, every $L\subseteq\F$ with
$|L|=n$, every $V=\RS[\E,L,d]$ of rate $\rho=d/n$, all words $u_0,\dots,u_l:L\to\E$, every
weight $\mu$ whose values have denominator $M$, every $\alpha\ge\alpha_*$, and every choice of
witnesses $P_z\in V$ with $\agree_\mu\big(\sum_{j\le l}z^ju_j,P_z\big)\ge\alpha$ for the
elements $z$ of a set $S\subseteq\F$: if $|S|>\Theta(n,M,l)$, then there are $v_0,\dots,v_l\in V$
with
\[
\mu\big(\{x\in L:\ u_j(x)=v_j(x)\ \forall j\}\big)\ge\alpha
\quad\text{and}\quad
\textstyle\sum_{j\le l}z^jv_j=P_z\ \text{ for more than }l\text{ values }z\in S .
\]
\end{hypothesis}
\uses{\cref{ext:bciks72,ext:proof-facts}}\usedby{\cref{prop:H1,prop:H2,cor:share,lem:round,thm:reduction}}

{\renewcommand{\theproposition}{P-H1}
\begin{proposition}[H holds at the Johnson radius]\label{prop:H1}
Let $\kappa\ge3$ and $\alpha_0:=\sqrt\rho\,(1+\frac1{2\kappa})$. Then
$\mathrm H(\alpha_0,\Theta_J)$ holds for
\[
\Theta_J(n,M,l):=l\,\max\Big(\frac{(1+\frac1{2\kappa})^7\kappa^7}{3\rho^{3/2}}\,n^2,\ \
\frac{2\kappa+1}{\sqrt\rho}\,(Mn+1)\Big).
\]
\end{proposition}

\begin{proof}
Let $\E,L,d,V,u_j,\mu,M,\alpha\ge\alpha_0,S,P_z$ be as in \ref{hyp:ca}, with
$|S|>\Theta_J(n,M,l)$. The proof of \cref{ext:bciks72} uses the set $S$ only through $|S|$ and
membership in it, and both places where the field enters hold over $\E$, since
$|S|\le|\F|\le|\E|$ and $\E$ has the characteristic of $\F$; so its conclusion holds for the
field $\E$ and the set $S$, giving $v_0,\dots,v_l\in V$ with the stated agreement. By the first
fact of \cref{ext:proof-facts} these $v_j$ satisfy $\sum_jz^jv_j=P_z$ for the given witnesses on
a set of more than $l$ values of $z$.
\end{proof}
\uses{\cref{ext:bciks72,ext:proof-facts}}\usedby{\cref{cor:johnson}}}

{\renewcommand{\theproposition}{P-H2}
\begin{proposition}[H holds in the unique-decoding regime]\label{prop:H2}
Put $\eta_{\mathrm U}:=\min\big(\frac{(1-\sqrt\rho)^2}2,\frac{\sqrt\rho}{20}\big)$. For every
$\alpha_*>\frac{1+\rho}2$, $\mathrm H(\alpha_*,\Theta_U)$ holds for
\[
\Theta_U(n,M,l):=l\,(Mn+1)\Big(\frac1{\eta_{\mathrm U}}+\frac3{\sqrt\rho}\Big).
\]
\end{proposition}

\begin{proof}
Let $\alpha\ge\alpha_*>\frac{1+\rho}2$ and $|S|>\Theta_U(n,M,l)$, and put
$\eta:=\min(\alpha-\sqrt\rho,\frac{\sqrt\rho}{20})$. Since
$\alpha-\sqrt\rho>\frac{1+\rho}2-\sqrt\rho=\frac{(1-\sqrt\rho)^2}2$ we have
$\eta\ge\eta_{\mathrm U}$, so $|S|>l(Mn+1)(\frac1\eta+\frac3{\sqrt\rho})$; also
$\alpha>\frac{1+\rho}2$ gives $\varepsilon=n/|\F|$ in $(\star\star)$, and
$(Mn+1)(\frac1\eta+\frac3{\sqrt\rho})>n$. So \cref{ext:bciks71} applies, with the same
transposition to $\E$ as in \cref{prop:H1}, and \cref{ext:facts71} supplies the curve identity
for the chosen witnesses on more than $l$ values of $z$.
\end{proof}
\uses{\cref{ext:bciks71,ext:facts71}}\usedby{\cref{cor:ud}}}
